% GENERATED FILE. Edit canonical lessons, then run npm run generate:books.
% canonical-source-hash: fnv1a64:8fb35413e9613198
% canonical-lessons: GU-C435-kadako, GU-C435-cakravat, GU-C435-susvato, GU-C435-agahi, GU-C435-jhabakvu

\chapter{Thunder and Cyclones}
\label{ch:gu-thunder-and-cyclones}
\hlchaptermodality{435}

\begin{chapteropening}
\textbf{By the end of this chapter:} \emph{I can talk about a storm and the forecast.}

Complete the last lesson of chapter 435: I can talk about a storm and the forecast.
\end{chapteropening}

\section[\texorpdfstring{\gu{કડાકો} (kaḍāko)}{કડાકો (kaḍāko)}]{\gu{કડાકો} (kaḍāko) — a crash of thunder; a loud crack}

\label{lesson:GU-C435-kadako}



\begin{quote}
Before the new one: say the Gujarati for to gather (of clouds), then the Gujarati for to get wet through.
\end{quote}

\subsection*{You'll want to know: \gu{કડાકો}}
\textbf{\gu{કડાકો}} — \emph{kaḍāko} — \textquotedblleft{}a crash of thunder; a loud crack\textquotedblright{}.

\begin{quote}
\gu{રાત્રે} \gu{વીજળીનો} \gu{મોટો} \gu{કડાકો} \gu{થયો}.
\end{quote}

\begin{quote}
\emph{rātre vījḷīno moṭo kaḍāko thayo.}
\end{quote}

\begin{quote}
At night there was a great crash of thunder.
\end{quote}

A masculine noun.

\begin{cousinweb}
\textbf{\gu{કડાકો}} (\emph{kaḍāko}) goes back to Sanskrit \emph{kaṭakaṭā}, an imitative word for a grating or crackling sound; CDIAL lists the Gujarati word, \textquotedblleft{}a loud crack\textquotedblright{} (CDIAL 2632).
\end{cousinweb}

\begin{grammarlens}[title={What you've built}]
One more word for getting around.
\end{grammarlens}

\subsection*{Your turn}
\begin{itemize}
\raggedright
  \item \emph{Say it:} \emph{kaḍāko}
  \item \emph{Say it:} \emph{kaḍāko}, once more
  \item \emph{Say it:} \emph{rātre vījḷīno moṭo kaḍāko thayo.} — At night there was a great crash of thunder.
  \item \emph{Recall:} say the Gujarati for to gather (of clouds), then the Gujarati for to get wet through, then say \emph{kaḍāko} again
\end{itemize}

\begin{tcolorbox}[breakable,colback=teal!4,colframe=teal!35!black,title={Before you move on}]
What does \gu{કડાકો} mean? (\textquotedblleft{}A crash of thunder; a loud crack\textquotedblright{}.) Say it once more.
\end{tcolorbox}
\section[\texorpdfstring{\gu{ચક્રવાત} (cakravāt)}{ચક્રવાત (cakravāt)}]{\gu{ચક્રવાત} (cakravāt) — a whirlwind; a cyclone (in weather reports)}

\label{lesson:GU-C435-cakravat}



\begin{quote}
Before the new one: say the Gujarati for to get wet through, then the Gujarati for a crash of thunder.
\end{quote}

\subsection*{You'll want to know: \gu{ચક્રવાત}}
\textbf{\gu{ચક્રવાત}} — \emph{cakravāt} — \textquotedblleft{}a whirlwind; a cyclone (in weather reports)\textquotedblright{}.

\begin{quote}
\gu{દરિયામાં} \gu{ચક્રવાત} \gu{આવ્યો}, \gu{એટલે} \gu{હોડીઓ} \gu{કિનારે} \gu{રહી}.
\end{quote}

\begin{quote}
\emph{dariyāmā̃ cakravāt āvyo, eṭle hoḍīo kināre rahī.}
\end{quote}

\begin{quote}
A cyclone came over the sea, so the boats stayed on the shore.
\end{quote}

A masculine noun.

\begin{cousinweb}
\textbf{\gu{ચક્રવાત}} (\emph{cakravāt}) is a tatsama, Sanskrit \emph{cakravāta}, \textquotedblleft{}a whirlwind\textquotedblright{} (Monier-Williams), made of \emph{cakra}, \textquotedblleft{}a wheel\textquotedblright{}, and \emph{vāta}, \textquotedblleft{}wind\textquotedblright{}: the wind that turns like a wheel.
\end{cousinweb}

\begin{grammarlens}[title={What you've built}]
One more word for getting around.
\end{grammarlens}

\subsection*{Your turn}
\begin{itemize}
\raggedright
  \item \emph{Say it:} \emph{cakravāt}
  \item \emph{Say it:} \emph{cakravāt}, once more
  \item \emph{Say it:} \emph{dariyāmā̃ cakravāt āvyo, eṭle hoḍīo kināre rahī.} — A cyclone came over the sea, so the boats stayed on the shore.
  \item \emph{Recall:} say the Gujarati for to get wet through, then the Gujarati for a crash of thunder, then say \emph{cakravāt} again
\end{itemize}

\begin{tcolorbox}[breakable,colback=teal!4,colframe=teal!35!black,title={Before you move on}]
What does \gu{ચક્રવાત} mean? (\textquotedblleft{}A whirlwind; a cyclone (in weather reports)\textquotedblright{}.) Say it once more.
\end{tcolorbox}
\section[\texorpdfstring{\gu{સુસવાટો} (susvāṭo)}{સુસવાટો (susvāṭo)}]{\gu{સુસવાટો} (susvāṭo) — a howl of wind; a whistling gust}

\label{lesson:GU-C435-susvato}



\begin{quote}
Before the new one: say the Gujarati for a crash of thunder, then the Gujarati for a whirlwind.
\end{quote}

\subsection*{You'll want to know: \gu{સુસવાટો}}
\textbf{\gu{સુસવાટો}} — \emph{susvāṭo} — \textquotedblleft{}a howl of wind; a whistling gust\textquotedblright{}.

\begin{quote}
\gu{રાત્રે} \gu{પવનના} \gu{સુસવાટા} \gu{સંભળાતા} \gu{હતા}.
\end{quote}

\begin{quote}
\emph{rātre pavannā susvāṭā sambhḷātā hatā.}
\end{quote}

\begin{quote}
At night you could hear the wind howling.
\end{quote}

A masculine noun.

\begin{cousinweb}
\textbf{\gu{સુસવાટો}} (\emph{susvāṭo}) goes back to \emph{*śūṣa}, \textquotedblleft{}blowing\textquotedblright{}, a form CDIAL reconstructs; CDIAL lists the Gujarati word, \textquotedblleft{}a blast of wind\textquotedblright{} (CDIAL 12577).
\end{cousinweb}

\begin{grammarlens}[title={What you've built}]
One more word for getting around.
\end{grammarlens}

\subsection*{Your turn}
\begin{itemize}
\raggedright
  \item \emph{Say it:} \emph{susvāṭo}
  \item \emph{Say it:} \emph{susvāṭo}, once more
  \item \emph{Say it:} \emph{rātre pavannā susvāṭā sambhḷātā hatā.} — At night you could hear the wind howling.
  \item \emph{Recall:} say the Gujarati for a crash of thunder, then the Gujarati for a whirlwind, then say \emph{susvāṭo} again
\end{itemize}

\begin{tcolorbox}[breakable,colback=teal!4,colframe=teal!35!black,title={Before you move on}]
What does \gu{સુસવાટો} mean? (\textquotedblleft{}A howl of wind; a whistling gust\textquotedblright{}.) Say it once more.
\end{tcolorbox}
\section[\texorpdfstring{\gu{આગાહી} (āgāhī)}{આગાહી (āgāhī)}]{\gu{આગાહી} (āgāhī) — a forecast}

\label{lesson:GU-C435-agahi}



\begin{quote}
Before the new one: say the Gujarati for a whirlwind, then the Gujarati for a howl of wind.
\end{quote}

\subsection*{You'll want to know: \gu{આગાહી}}
\textbf{\gu{આગાહી}} — \emph{āgāhī} — \textquotedblleft{}a forecast\textquotedblright{}.

\begin{quote}
\gu{હવામાન} \gu{ખાતાએ} \gu{ભારે} \gu{વરસાદની} \gu{આગાહી} \gu{કરી}.
\end{quote}

\begin{quote}
\emph{havāmān khātāe bhāre varsādnī āgāhī karī.}
\end{quote}

\begin{quote}
The weather office forecast heavy rain.
\end{quote}

A feminine noun.

\begin{cousinweb}
\textbf{\gu{આગાહી}} (\emph{āgāhī}) is usually traced to Persian \emph{āgāhī}, \textquotedblleft{}awareness, being told in advance\textquotedblright{}, built on \emph{āgāh}, \textquotedblleft{}aware\textquotedblright{}.
\end{cousinweb}

\begin{grammarlens}[title={What you've built}]
One more word for getting around.
\end{grammarlens}

\subsection*{Your turn}
\begin{itemize}
\raggedright
  \item \emph{Say it:} \emph{āgāhī}
  \item \emph{Say it:} \emph{āgāhī}, once more
  \item \emph{Say it:} \emph{havāmān khātāe bhāre varsādnī āgāhī karī.} — The weather office forecast heavy rain.
  \item \emph{Recall:} say the Gujarati for a whirlwind, then the Gujarati for a howl of wind, then say \emph{āgāhī} again
\end{itemize}

\begin{tcolorbox}[breakable,colback=teal!4,colframe=teal!35!black,title={Before you move on}]
What does \gu{આગાહી} mean? (\textquotedblleft{}A forecast\textquotedblright{}.) Say it once more.
\end{tcolorbox}
\section[\texorpdfstring{\gu{ઝબકવું} (jhabakvũ)}{ઝબકવું (jhabakvũ)}]{\gu{ઝબકવું} (jhabakvũ) — to flash (of lightning); to flicker}

\label{lesson:GU-C435-jhabakvu}



\begin{quote}
Before the new one: say the Gujarati for a howl of wind, then the Gujarati for a forecast.
\end{quote}

\subsection*{You'll want to know: \gu{ઝબકવું}}
\textbf{\gu{ઝબકવું}} — \emph{jhabakvũ} — \textquotedblleft{}to flash (of lightning); to flicker\textquotedblright{}.

\begin{quote}
\gu{રાત્રે} \gu{આકાશમાં} \gu{વીજળી} \gu{ઝબકી}.
\end{quote}

\begin{quote}
\emph{rātre ākāśmā̃ vījḷī jhabkī.}
\end{quote}

\begin{quote}
At night lightning flashed in the sky.
\end{quote}

Past (it …): \gu{તે} \gu{વીજળી} \gu{ઝબકી}, \emph{te vījḷī jhabkī}.

\begin{cousinweb}
\textbf{\gu{ઝબકવું}} (\emph{jhabakvũ}) goes back to \emph{*jhamm}, \textquotedblleft{}a flash\textquotedblright{}, a form CDIAL reconstructs; CDIAL lists the Gujarati verb, \textquotedblleft{}to flash\textquotedblright{} (CDIAL 5341).
\end{cousinweb}

\begin{grammarlens}[title={What you've built}]
One more word for getting around.
\end{grammarlens}

\subsection*{Your turn}
\begin{itemize}
\raggedright
  \item \emph{Say it:} \emph{jhabakvũ}
  \item \emph{Say it:} \emph{jhabakvũ}, once more
  \item \emph{Say it:} \emph{rātre ākāśmā̃ vījḷī jhabkī.} — At night lightning flashed in the sky.
  \item \emph{Recall:} say the Gujarati for a howl of wind, then the Gujarati for a forecast, then say \emph{jhabakvũ} again
\end{itemize}

\begin{tcolorbox}[breakable,colback=teal!4,colframe=teal!35!black,title={Before you move on}]
What does \gu{ઝબકવું} mean? (\textquotedblleft{}To flash (of lightning); to flicker\textquotedblright{}.) Say it once more.
\end{tcolorbox}
